\documentclass[11pt]{article}
\usepackage[margin=1.1in]{geometry}
\usepackage{amsmath,amssymb,amsthm}
\usepackage{enumitem}
\usepackage[colorlinks=true,linkcolor=blue,urlcolor=blue]{hyperref}

\newcommand{\Stfa}{\mathrm{St}_{\mathrm{fa}}}
\newcommand{\Stsig}{\mathrm{St}_{\sigma}}
\newcommand{\HAND}{\ensuremath{[\![\text{HAND}]\!]}}
\newcommand{\file}[1]{\texttt{\small #1}}
\newcommand{\gap}[1]{\medskip\noindent\textbf{Gap #1.}}
\newcommand{\lookfor}{\par\smallskip\noindent\emph{What to look for:} }
\newcommand{\redflag}{\par\smallskip\noindent\emph{Red flag:} }

\title{Ratification kit: the hand-verification gaps\\
\large What the scripts could not check, item by item, with references}
\author{Prepared end of session 16}
\date{2026-07-11}

\begin{document}
\maketitle

\noindent\textbf{How to use this document.} Every item below was
verified by a fresh-context adversarial review with from-scratch
machine receipts; each receipt honestly lists the steps that are
\emph{hand-proof-only} --- infinite objects, topology, choice
principles, and source-reading judgment calls that no finite script
reaches. Those steps are where your eye is the only remaining
verification layer. For each item: the claim chain, the receipt, then
each gap with (a) exactly what to check, (b) where the written proof
lives, (c) the primary reference with page, (d) what a failure would
look like and what it would take down. Ratifying an item = you accept
its gaps after inspection (or spot-check the ones you choose and
accept the receipt for the rest); tell me and I mark the records.

\smallskip
\noindent\textbf{Suggested order.} Items 2--3--4 are serial (each uses
the previous); item 1 is quick and independent; item 5 is a separate
lane entirely. Item 2's gaps are the highest-leverage in the corpus:
A2 and T3 carry everything after them.

\tableofcontents

%=====================================================================
\section{Item 1 (s9): the MPT theorem-let --- no non-Boolean concrete
$\sigma$-class OML has all $\sigma$-additive states Jauch--Piron}

\noindent\textbf{Claim chain} (attack note \S7c(i)):
on a concrete $\sigma$-class, every Dirac state is two-valued and
$\sigma$-additive [glue, \HAND]; all-Dirac-states-JP
$\Leftrightarrow$ $C_{\mathrm{card}}$ [MPT Prop 3.2, \emph{their}
proof]; $C_{\mathrm{card}} \Rightarrow$ downward-directed [MPT p.~10,
one direction only]; downward-directed $+$ lattice $\Rightarrow$
Boolean [MPT Prop 4.4, \emph{their} proof]. Consequences: BNPW 1985
excluded from $\mathrm{Adm}\cap\mathrm{OML}$ under every reading
(disjunction); Skeleton C first half refuted by MO$_2$.

\noindent\textbf{Receipt:} \S7c(iii-b)/(iii-c) of
\file{../oml\_attack/oml\_lattice\_regularity\_attack.md};
\file{../verification/mo2\_jp\_check.py} (seven properties). Reviewer
read the MPT PDF in full.

\smallskip\noindent\textbf{Lean note (s17--s18):} gap 1.1's glue has
been in Lean since the spine: \texttt{SigmaEssentialLocalization.dirac}
proves $\delta_\omega$ is a \texttt{TwoValuedState} (two-valued,
$\sigma$-additive) on ANY concrete $\sigma$-class --- the
\texttt{val\_iUnion} field is exactly ``$x$ lies in exactly one
$E_n$.'' Only the scope-line reading (Adm's completeness sense) and
gaps 1.2--1.3 (source-reading) remain for the eye.

\gap{1.1} \emph{The glue: Dirac states are $\sigma$-additive on
$\sigma$-classes.} Three lines: if $E = \biguplus_n E_n \in L$
(disjoint union), then $x$ lies in exactly one $E_n$, so
$\delta_x(E) = \sum_n \delta_x(E_n)$.
\lookfor the only subtlety is that $\sigma$-class
$\sigma$-completeness is the \emph{union} sense (disjoint countable
unions are set unions and lie in $L$), matching
\file{papers/sigma\_essential/sigma\_essential\_body.tex} q:bare clause (b). Confirm the
theorem-let's scope line: ``Adm's $\sigma$-completeness IS the
$\sigma$-class sense.''
\redflag if Adm's completeness were \emph{sup}-based (suprema exist
but are not set unions), the glue fails --- that is exactly the
distinction MPT Prop 4.6 (p.~11) navigates, and the reviewer checked
the senses match. Reference: M\"uller--Pt\'ak--Tkadlec, \emph{Int.\ J.
Theor.\ Phys.} 31 (1992) 843--854, PDF
\file{muller\_ptak\_tkadlec\_1992\_covering\_properties.pdf}: Def
1.1 (p.~2, \emph{finite} closure only --- their default logic is
weaker than a $\sigma$-class), Prop 3.2 (p.~6, no side conditions),
Prop 4.4 (p.~10).

\gap{1.2} \emph{Direction of the Prop 3.2 bridge.} Prop 3.2's
$C_{\mathrm{card}}$ is per-point ($\forall x \in A\cap B\ \exists C
\ni x$, $C \subseteq A \cap B$), which implies downward-directedness
strictly --- it is NOT an iff.
\lookfor \S7c(i) uses only the correct direction
($\Rightarrow$); check no later use silently upgrades it.

\gap{1.3} \emph{The BNPW disjunction} (\S7c, ``Consequence for
BNPW''). Pure logic over an underdetermined source: either BNPW's
object is a $\sigma$-class (then not a lattice, by the theorem-let)
or it is not $\sigma$-closed (then not an Adm-style carrier).
\lookfor the disjunction exhausts the readings of MPT's Q5.4 note
(p.~12); the Bunce--Hamhalter 2000 ILL pull remains
confirmation-grade only --- nothing is cited from it.
\redflag any future text citing BNPW as a lattice witness or
non-witness \emph{simpliciter}.

%=====================================================================
\section{Item 2 (s11 math + s12 edits): attack note \S9 --- the
foundation layer (A1, L0, A2, T1, T3, P1, Conjecture $B'$)}

\noindent\textbf{Why this item matters most:} A2 (blocks are
$\sigma$-fields) and T3 (Dirac realization on countably generated
blocks) are load-bearing for every subsequent session (\S10's
factoring, \S11's $P^=$ and ladder). A defect here propagates
everywhere.

\noindent\textbf{Receipt:} \file{../oml\_attack/PROOF\_READ\_2026-07-10\_attack\_s9.md}
(118/118 checks, three flagged high-risk steps cleared);
scripts \file{../verification/proof\_read\_2026-07-10\_s11/}.

\smallskip\noindent\textbf{Lean bolster (2026-07-11, s17--s18 --- added
after this kit was issued).}
\file{formalization/QuerySystem/QuerySystem/ConcreteOMLBlocks.lean} and
\file{.../MarczewskiTransport.lean} (build; \texttt{\#print axioms}
receipts at file end) convert this item to machine receipts, with
\textbf{NO axioms anywhere} --- classical base only:
\begin{itemize}[leftmargin=2em]
\item \textbf{Gap 2.1 DISCHARGED} --- A2 in full
  (\texttt{IsMaxBlock.inter\_mem}, $\sigma$-field packaging
  \texttt{IsMaxBlock.toMeasurableSpace}, block overlaps
  \texttt{overlapMeasurableSpace}). The anticipated Foulis--Holland
  citation axiom proved UNNECESSARY: the needed instance follows from a
  three-line \emph{meet squeeze} (\texttt{compat\_of\_locally\_resolved}
  --- the meet of $X, S$ contains every carrier subset of $X \cap S$, so
  a pointwise carrier cover forces meet $= X \cap S \in L$), applied
  with the cover $\{A \cap B, A \cap C\}$ of $A \cap (B \cup C)$; set
  distributivity supplies what orthomodular calculus supplies
  abstractly. Commutant closure = \texttt{compat\_inter\_left}, proved.
  Bruns--Harding remains the citation for the ABSTRACT theorem only;
  nothing in the corpus's chain now rests on it.
\item \textbf{Gap 2.2 DISCHARGED} --- T3 in full
  (\texttt{dirac\_realization\_of\_countablyGenerated}), including
  $\sigma$-continuity from above and the dichotomy $\sigma$-subfield;
  stated on abstract countably generated $\sigma$-fields, so no
  circularity with A2.
\item \textbf{Gap 2.3(b) DISCHARGED} --- the $\omega_1$ counterexample
  is machine-checked (\texttt{MarczewskiTransport.lean}): the
  co-countable state is a $\sigma$-additive \texttt{TwoValuedState}
  (\texttt{cocountableState}), non-Dirac
  (\texttt{cocountableState\_not\_isDirac}), and the field is NOT
  countably generated (\texttt{cocountable\_not\_countably\_generated}
  --- itself a corollary of T3). Gap 2.3(a) (product-Ulam pattern
  intersections escape $L$) stays with the oracle scripts ---
  corroboration only.
\item \textbf{Gaps 2.4, 2.5 DISCHARGED} --- A1(a)--(c)
  (\texttt{exists\_isMaxBlock\_superset}, \texttt{IsMaxBlock.iUnion\_mem},
  \texttt{isSigmaOn\_carrier\_iff\_maxBlocks}); L0 both directions
  (\texttt{Compat.isGreatest\_inter},
  \texttt{compat\_of\_commuting\_decomp}). Bonus: T1 in both forms
  (\texttt{compat\_of\_singletons}, \texttt{interClosed\_of\_singletons})
  and P1's poor-pair anatomy (\texttt{PoorPair.*}) are also machine
  receipts now.
\end{itemize}
Every receipt in the chain lists exactly \texttt{[propext,
Classical.choice, Quot.sound]}. Ratifying item 2 = accepting the build.

\gap{2.1} \emph{A2's infinite commutant case.} The step ``every
pairwise-compatible subset --- of any cardinality --- generates a
Boolean subalgebra'' rests on Bruns--Harding Prop 2.8, checked by the
scout to be stated for \emph{arbitrary} subsets (via commutant
closure, not a directed-union limit).
\lookfor read Bruns--Harding 2000, ``Algebraic Aspects of Orthomodular
Lattices'' (in \emph{Current Research in Operational Quantum Logic},
Kluwer), Props 2.2--2.4 (commutant $C(a)$ closed under $'$ and all
existing joins/meets), Prop 2.8 (arbitrary pairwise-commuting subsets),
Prop 3.1 (blocks). Verify Prop 2.8's statement covers infinite sets
with no separability hypothesis. Nuance already pinned in the receipt:
Foulis--Holland (Kalmbach 1983, Thm 5 p.~25; Holland 1964 \emph{TAMS}
112, \textbf{not} 1963) gives only distributivity of the generated
sub\emph{lattice}; the Boolean sub\emph{algebra} conclusion for
infinite families is exactly what Prop 2.8 adds.
\redflag if Prop 2.8 needed a finiteness or completeness hypothesis,
A2's ``blocks are $\sigma$-fields'' would lose the infinite case, and
with it T3, $P^=$, and the \S10/\S11 superstructure.

\gap{2.2} \emph{T3's $\sigma$-continuity derivation} (infinite; the
scripts only exercised finite $\sigma$-fields and the pentagon).
Re-derive: $F_n = $ intersection of the first $n$ generators'
value-1 sides ($\in Bl$ by A2); $\nu(F_n) = 1$ by two-valued f.a.\
induction; $D = \bigcap_n F_n \in Bl$;
$F_1 = D \uplus \biguplus_n (F_n \setminus F_{n+1})$ with each
difference $\nu$-null, so $\sigma$-additivity \emph{on} $Bl$
(supplied by A1c) gives $\nu(D) = 1$, hence $D \neq \emptyset$;
finally $\{A \in Bl : D \subseteq A \text{ or } D \cap A =
\emptyset\}$ is a $\sigma$-subfield containing the generators.
\lookfor two hidden joints: (a) ``countable unions in $Bl$ ARE set
unions'' --- that is A2, so check there is no circularity (A2's proof
uses only Bruns--Harding + L0 + A1b, never T3 --- confirm by reading
\S9c's A2 paragraph); (b) the $\sigma$-subfield step needs closure
under countable unions of the dichotomy family, one line but write it.
\redflag a use of $\sigma$-continuity \emph{from above} on $L$ itself
(not available in a $\sigma$-class); the proof must route through the
block, where A2 makes it a genuine $\sigma$-field.

\gap{2.3} \emph{T3's load-bearing counterexamples.} (a) Without
latticehood the $F_n$ need not lie in $L$: this is the corpus's own
Cor.\ ``incompatibility; not a lattice''
(\file{papers/sigma\_essential/sigma\_essential\_body.tex:713}) --- the reviewer marked it
consistent-but-not-re-derived (receipt finding 6).
\lookfor open that Corollary once and confirm it states what T3 needs:
on the product Ulam carrier the finite intersections of pattern sets
escape $L$. (b) Without countable generation: the
countable/co-countable field on $\omega_1$ with the co-countable
state --- re-derive that the state is $\sigma$-additive (a countable
disjoint union of countable sets is countable, so the co-countable
member is unique) and that the field is genuinely not countably
generated (any countable family of countable/co-countable sets
generates a field whose atoms include an uncountable set that splits
further).

\gap{2.4} \emph{A1(b): blocks closed under countable disjoint unions.}
The orthomodular-difference step: for $u = \biguplus c_n$, $c_n \in
Bl$, and $b \in Bl$: $u \cap b = \biguplus (c_n \cap b) \in L$, and
$u \setminus b,\ b \setminus u \in L$ likewise, so $u$ is compatible
with all of $Bl$, hence in $Bl$ by maximality.
\lookfor each $c_n \cap b \in L$ needs $c_n \leftrightarrow b$ (both
in $Bl$) $+$ L0; the disjointness of $(c_n \cap b)_n$ is inherited;
maximality is the Zorn-supplied property from A1(a).

\gap{2.5} \emph{L0's converse (join vs.\ $\uplus$).} $P \uplus Q \in
L$ is an upper bound so $P \vee Q \subseteq P \uplus Q$; the join
contains $P \cup Q$ as a set since it bounds each part; equality
follows, then $A \cap B = A \wedge B$.
\lookfor this is the step the s12 reviewer singled out and closed;
re-run it once --- it is four lines and everything concrete-vs-lattice
downstream leans on it.

\smallskip\noindent\textbf{s12 edits to sign off} (all cosmetic, listed
in the receipt): pentagon irreducibility parenthetical; T1 hypothesis
cleaned (every point of $A \cap B$ lies in some member of $L$ inside
it --- \emph{no exhaustion needed}); corollary remark replaced;
P1's use of poorness made explicit; $B'$(ii) precision caveat
(whole-carrier property phrased as coarse-block --- must be sharpened
before attack).

%=====================================================================
\section{Item 3 (s13 math + s14 edits): attack note \S10 --- Skeleton
A read, theorem-lets R and P, the demotion and factoring}

\noindent\textbf{Claim chain:} R ((8.1)-inner-regular two-valued f.a.\
state on Hausdorff $=$ Dirac restriction with compact witness); P
(diffuse countable/co-countable components fail (8.1) on uncountable
Polish); consequence C (coarse-riding states escape DW Thm D.6 on
\emph{every} Polish representation) $\Rightarrow$ Skeleton A demoted,
Theorem 2 factors as $B'$(i) $+$ coarse-block question.

\noindent\textbf{Receipt:} \file{../oml\_attack/PROOF\_READ\_2026-07-10\_attack\_s10.md}
(11{,}686/11{,}686 checks; sources re-checked verbatim);
scripts \file{../verification/proof\_read\_2026-07-10\_s13/}.

\smallskip\noindent\textbf{Lean bolster (s18):}
\file{formalization/QuerySystem/QuerySystem/InnerRegularity.lean}
(builds; receipts axiom-free):
\begin{itemize}[leftmargin=2em]
\item \textbf{Gap 3.1 DISCHARGED} --- theorem-let R in full
  (\texttt{dirac\_forcing}): honest field-of-sets hypotheses
  (\texttt{SetField}/\texttt{FieldState}, finite closure only, no
  $\sigma$ anywhere), $\mathcal K_1 \neq \emptyset$ and $D \neq
  \emptyset$ as CONCLUSIONS (the ✎s14 rewrite), the FIP/finite-subcover
  step machine-checked, Hausdorff consumed exactly at
  compact-$\Rightarrow$-closed. The red-flag counterexample is ALSO
  machine-checked: on \texttt{CofiniteTopology} every set is compact
  (\texttt{isCompact\_cofiniteTopology}), (8.1) holds
  (\texttt{cofinite\_81\_holds}), the cofinite state is non-Dirac ---
  and R itself then yields \texttt{not\_t2Space\_cofiniteTopology}.
\item \textbf{Gap 3.2 DISCHARGED (two-valued instance, strengthened)}
  --- \texttt{cocountable\_fails\_81\_at}: on any uncountable
  SECOND-COUNTABLE Hausdorff space (subsuming Polish), the co-countable
  state fails (8.1) at the specific $F = X \setminus \{x\}$. The
  topological content is machine-checked via condensation points
  (\texttt{exists\_condensation\_point}), replacing the
  Cantor--Bendixson route entirely; ``closed co-countable $\supseteq$
  perfect kernel'' is consumed only at the chosen point. The
  real-valued diffuse-mass decomposition stays hand.
\item \textbf{Gap 3.6 DISCHARGED (two-valued instance)} ---
  \texttt{cocountable\_fails\_81\_hausdorff}: on ANY uncountable
  Hausdorff space, failure somewhere by reduction to R.
\item Gaps 3.3 (witnessed moot), 3.4 (DW fn-29), 3.5 (Maharam anatomy)
  are source-reading --- unchanged, yours.
\end{itemize}

\gap{3.1} \emph{R's FIP step for infinite $\mathcal K_1$.} The
value-1 in-field compacts have the FIP \emph{inside one fixed compact
member}; nonempty intersection needs the members closed in that
compact (Hausdorff: compact $\Rightarrow$ closed).
\lookfor the s14 rewrite of R's statement (asserting $\mathcal K_1
\neq \emptyset$ and $D \neq \emptyset$ as conclusions, avoiding the
$\bigcap \emptyset = X$ misparse); then check Hausdorff is genuinely
load-bearing via the receipt's counterexample: $\mathbb N$ with the
cofinite topology, finite/cofinite field, cofinite state ---
everything is compact, (8.1) trivially holds, no Dirac.
\redflag any restatement of R that drops Hausdorff or quantifies
(8.1) over topological compacts not sandwiched by in-field sets
(that is exactly the DW fn-29 gap; see 3.4).

\gap{3.2} \emph{P's topological content} (not scriptable):
Cantor--Bendixson on an uncountable Polish space; the perfect kernel
is uncountable; compact $\Rightarrow$ closed; and the sub-claim
``closed co-countable $\supseteq$ perfect kernel.''
\lookfor re-derive the sub-claim: if $F$ is closed and co-countable,
$X \setminus F$ is open and countable, hence contains no point of the
perfect kernel (a nonempty open subset of a perfect Polish subspace is
uncountable); the receipt's second counterexample shows ``$x \in$
perfect kernel'' is load-bearing ($X = [0,1] \uplus \{p\}$ isolated,
$x = p$ breaks the argument). Reference for background: any DST text
(Kechris \S6 for Cantor--Bendixson) --- also on the skill plan.

\gap{3.3} \emph{P0's choice flag.} Atom-set countability uses
$\mathrm{AC}_\omega(\mathrm{fin})$ on abstract $X$; eliminable on
Polish $X$ (definable order).
\lookfor confirm P0 is only ever \emph{applied} on Polish carriers in
\S10 (it is), so the flag is moot in context; nothing to fix, just
witness it.

\gap{3.4} \emph{The DW footnote-29 reading.} The claim ``DW's
blockwise inner-regularity leg omits $K \in \mathcal A_i$, gap
confirmed in the text, no repair in Appendix D'' is a
\emph{source-reading} judgment.
\lookfor Derr--Williamson 2023, arXiv:2302.03522
(\file{derr\_williamson\_2023.pdf}): Thm D.6 p.~48 (five legs;
blockwise inner regularity is a distinct leg), footnote 29 (verbatim,
membership-free), proof p.~49 (``all restrictions $\mu_i$ are inner
regular'' --- no repair). Then check the corpus convention: we cite
D.6 only in the (8.1) sense (Maharam-style in-field sandwich), which
is the reading that makes their proof work.
\redflag citing D.6 with fn-29's literal (membership-free) reading
anywhere --- the s13 Addendum exists precisely to pin this.

\gap{3.5} \emph{Maharam anatomy M1--M3} (source-reading; verbatim
anchors already in the receipt).
\lookfor Maharam 1972 (\file{maharam\_1972.pdf}): (8.1) display
p.~145 (with $K \in \mathcal F_\alpha$ --- the in-field membership DW's
footnote drops); the Hahn--Banach engine is her \emph{Lemma} 6.1 (Thm
6.1 is the criterion invoking it); her Remark p.~146 (the $m_\alpha$
``have countably additive extensions''). Her proof of 8.1 (p.~146)
uses (8.1) only for the compact/open sandwich and the $K \subseteq
G_1 \cup \dots \cup G_n$ step --- this is what licenses \S10's
``(8.1) consumed ONLY as intersection-stable compact class.''

\gap{3.6} \emph{The all-Hausdorff strengthening} (banked s14): diffuse
mass $c > 0$ kills (8.1) on \emph{every} Hausdorff topology, by
reduction to R.
\lookfor the reduction: an in-field compact $K \subseteq F$ with
$\mu(K) > \mu(F) - c$ must be co-countable, so the co-countable
component $\nu$ would itself satisfy (8.1), contradicting R ($\nu$
non-principal). The threshold arithmetic is machine-checked; the
topological conclusion is the hand part.

\smallskip\noindent\textbf{s14 edit to sign off:} the scope fix on
``fails \emph{exactly} on the coarse locus'' --- earned only among
\emph{coarse-block $\sigma$-states} (atomic states satisfy (8.1) there
via finite in-field atom-truncations; compact-poor blocks fail (8.1)
even at Diracs: $\delta_\omega\!\restriction\!\{\emptyset,
\mathrm{irr}, \mathbb Q, \mathbb R\}$, $\omega$ irrational).
Propagated to \S10d(i) and the prior-art verdict banner. The s13
conclusions use only the failure direction and are unaffected ---
verify that one sentence in the receipt and you have verified the
propagation's safety.

%=====================================================================
\section{Item 4 (s15 math + s16 edits): attack note \S11 --- the
$B'$(i) reduction ladder and the Marczewski read}

\noindent\textbf{Claim chain} (as reviewed, C1--C12): 2BR two-block
rescue; cluster normal form; $P^=$ (pointed $\Leftrightarrow$
$\sigma$ on countably generated blocks); kernel-FIP strictly stronger
than coherence; representation-dependence of $K(s) = \emptyset$;
$\Phi \Leftrightarrow \Stsig$ dense in $\Stfa$; in-block 4(i) $+$
compact-transport criterion; $\omega_1$ scoping; no-cheap-union;
T4 ladder.

\noindent\textbf{Receipt:} \file{../oml\_attack/PROOF\_READ\_2026-07-11\_attack\_s11.md}
(424{,}509/424{,}509 checks); scripts
\file{../verification/proof\_read\_2026-07-11\_s15/}; Marczewski
source-check inside the receipt.

\smallskip\noindent\textbf{Lean bolster (s18):}
\file{formalization/QuerySystem/QuerySystem/ConcreteOMLPatterns.lean} +
\file{.../MarczewskiTransport.lean} (build; receipts axiom-free):
\begin{itemize}[leftmargin=2em]
\item \textbf{Gap 4.1 DISCHARGED} --- $P^=$ in full, split honestly:
  ($\Leftarrow$) \texttt{isSigmaOn\_of\_blockwisePointed} needs NO
  countable generation and NO latticehood (the banked sharpening);
  ($\Rightarrow$) \texttt{blockwisePointed\_of\_isSigmaOn} via T3;
  iff = \texttt{pointed\_iff\_sigma}; the selection reading =
  \texttt{phiCluster\_iff\_pointedSelection}. Also 2BR
  (\texttt{two\_block\_rescue}) and the cluster machinery
  (\texttt{phi\_iff\_phiCluster}, \texttt{cluster\_extension}) --- the
  general-$\Omega$ halves of gap 4.6.
\item \textbf{Gap 4.2 DISCHARGED} --- $\Phi$-density
  (\texttt{phi\_iff\_dense}) on the Cantor cube over the carrier;
  (a) \texttt{isClosed\_stFA} (the f.a.\ states = the closed constraint
  set, the onto/nonempty-traces reading built into the proof);
  (b) \texttt{isCompact\_stFA} (Tychonoff, ZFC).
\item \textbf{Gap 4.3 DISCHARGED} --- the $\omega_1$ scoping bank:
  \texttt{cocountableClass\_countablyCompact} (with $\emptyset$
  adjoined, the ✎s16 fix visible in the FIP-vacuous branch) and
  \texttt{cocountableClass\_approximates}.
\item \textbf{Gap 4.4(a) DISCHARGED} --- the CORRECTED
  compact-transport criterion is a theorem:
  \texttt{IsMaxBlock.isSigmaOn\_of\_compact\_class} (in-block) and
  \texttt{compact\_transport} (OML form, via A1c), with Marczewski
  approximation --- the value-1 inner witness $D \subseteq K \subseteq
  E$ --- as the hypothesis, exactly the ✎s16 form. The
  singleton-class inequivalence can no longer regress silently: the
  criterion's Lean hypothesis IS approximation. 4.4(b)/(c)
  (Marczewski attribution) stay source-reading.
\item \textbf{Gap 4.5, clean half DISCHARGED} --- \texttt{t4At\_of\_phi}
  / \texttt{bPrimeI\_implies\_T4}: B$'$(i) $\Rightarrow$ T4, with the
  T3 kernel proved to be an ATOM of the block. \texttt{BPrimeI} and
  \texttt{T4} are named open \texttt{Prop}s; the crux is NOT stated as
  a theorem anywhere --- the red flag is structurally respected.
\item Gap 4.6's remainders are covered en passant (monotonicity is
  \texttt{FinAddState.val\_mono}, lattice-free by construction).
\end{itemize}

\gap{4.1} \emph{$P^=$'s infinite legs} (scripts cover only the finite
pentagon). Three joints: (a) the ($\Rightarrow$) leg is T3 run on a
countable generating family --- inherits item 2's gap 2.2; (b) the
kernel-nonempty $\Rightarrow$ charges-an-atom bridge: for $\omega \in
D_{Bl}$, if $\nu(A(\omega)) = 0$ then $\nu(A(\omega)^c) = 1$ forces
$\omega \in A(\omega)^c$, absurd --- needs C4's atomicity; (c) the
($\Leftarrow$) leg: $\omega \in D_{Bl}$ makes
$\nu\!\restriction\!Bl = \delta_\omega\!\restriction\!Bl$ (two-valued
$+$ field), which is $\sigma$-additive per block, then A1c globalizes.
\lookfor C4's ``every member of $Bl$ is a union of signature cells'':
the family of cell-saturated sets is a $\sigma$-field containing the
generators --- write the countable-union line once. Note the banked
sharpening: ($\Leftarrow$) needs \emph{no} countable generation.
\redflag any later use of $P^=$ on a block \emph{not} countably
generated in the ($\Rightarrow$) direction --- that direction is
exactly what coarse blocks refuse ($\omega_1$ counterexample).

\gap{4.2} \emph{$\Phi$-density's topology} (C8). Three hand steps:
(a) $\Stfa \subseteq \{0,1\}^L$ is closed --- each finite-additivity
constraint involves finitely many coordinates, so violations are
witnessed on basic clopens; (b) compactness of $\{0,1\}^L$ is
Tychonoff (BPI-strength --- fine in ZFC, flagged for the record);
(c) the pattern-vs-clopen correspondence: a finite value assignment
with an f.a.\ witness $\perp\uplus$-closes to a finite
$\perp$-closed pattern \emph{without leaving the clopen} (the closure
only adds forced values). The map pattern $\mapsto$ trace is onto,
not injective --- harmless, but the ``correspond exactly'' phrasing
should be read as ``nonempty traces.''
\lookfor then the equivalence itself: $\Phi$ says every nonempty
basic clopen of $\Stfa$ meets $\Stsig$ --- literally density. Check
$\Phi$'s quantifier in \file{papers/sigma\_essential/sigma\_essential\_body.tex}
(\S Discussion) is over \emph{finite} patterns, which it is.

\gap{4.3} \emph{The $\omega_1$ scoping bank} (C10). (a) The
co-countable family (with $\emptyset$ adjoined --- s16 fix) is a
countably compact class: a countable intersection of co-countables is
co-countable, nonempty since $\omega_1$ is uncountable (uses
$\mathrm{AC}_\omega$: countable unions of countable sets); any
sequence containing $\emptyset$ vacuously fails the FIP hypothesis.
(b) It Marczewski-approximates the countable/co-countable field
w.r.t.\ the coarse state: co-countable $E$ takes $D = K = E$;
countable $E$ takes $K = \emptyset$. (c) Independent corroboration:
the state is purely atomic ($\omega_1$ one atom), hence compact by
M--RN \S3(ii) directly.
\lookfor Marczewski, ``On compact measures,'' \emph{Fund.\ Math.} 40
(1953) 113--124 (\file{marczewski\_1953\_on\_compact\_measures.pdf}):
compact class \S2 p.~115 (sequence form; purely set-theoretic --- his
stated aim p.~113 is ``eliminating non-essential topological
concepts''); approximation \S3 p.~116 ($\eta$-form over \emph{all}
$E \in \mathbf M$ --- this is why $\emptyset$ is needed); compact
measure \S4 p.~118; 4(i) p.~118.
\redflag treating Marczewski compactness as if it implied topological
(8.1)-regularity --- the whole point of the bank is that the coarse
state is Marczewski-compact yet kills the topological leg.

\gap{4.4} \emph{The s16 wording fixes --- these change standing
conjecture targets; sign off explicitly.} (a) Compact-transport
criterion and the coarse-factor wall statement now read
``Marczewski-approximates in-block'' instead of ``refines the value-1
filter.'' The two are \emph{inequivalent}: the class of all singletons
of $\Omega$ is countably compact and refines every ultrafilter, which
would make every two-valued f.a.\ state $\sigma$-additive ---
absurd (any nonprincipal ultrafilter state on $\mathcal P(\mathbb N)$).
The value-1 \emph{inner} witness $D \in Bl$, $\mu(D) = 1$,
$D \subseteq K$ is what powers the FIP step.
\lookfor re-run the criterion's proof once with the corrected
hypothesis: tails $F_n$ value-1, sandwiches $D_n \subseteq K_n
\subseteq F_n$, $\mu(D_n) = 1$; finite intersections of the $D_n$
are value-1 (in-block ultrafilter), hence nonempty, giving the $K_n$
the FIP; countable compactness yields $\emptyset \neq \bigcap K_n
\subseteq \bigcap F_n = \emptyset$. The coarse-block toy hunt will
aim at the wall statement in this corrected form --- this is the one
edit with forward-pointing consequences.
(b) 5(iii)/5(iv): the clean ``independent $\sigma$-fields $+$ compact
partials $\Rightarrow$ compact'' is 5(iv) (p.~120); 5(iii) puts its
hypotheses on the \emph{approximating classes} and partial
compactness alone is explicitly insufficient there (``and, what is
more''). Both are compactness-\emph{transfer} theorems, not
extension-existence (that is Marczewski 1951, \emph{Fund.\ Math.} 38,
Thm I). (c) M--RN \S4 credit: the positive half (minimal
$\sigma$-extensions of compact measures are compact) is M 1953 4(ii),
quoted in M--RN as ``C 4 (ii)''; M--RN \S4's own theorem is the
converse-false example. Reference:
\file{marczewski\_ryll\_nardzewski\_1953\_compactness\_products.pdf},
pp.~165--170.

\gap{4.5} \emph{The Zorn frame and the crux scope} (C12(iii) $+$ s16
scope note). T4's \emph{necessity} for $B'$(i) is clean ($B'$(i)
$\Rightarrow$ $\sigma$-extension of the cluster-pattern
$\Rightarrow$ pointed at every block by $P^=$ $\Rightarrow$ the
charged atom of $Bl_0$ witnesses coherence). But the boxed crux is at
\emph{finite} support, and the Zorn run needs the successor at
arbitrary $\mathcal S$: a block may have infinitely many atoms, so
``point $Bl_0$ somehow'' is an infinite union of closed conditions,
not closed --- compactness alone does not close the limit stage.
\lookfor the s16 annotation at the ``everything reduces'' sentence;
confirm you are happy with the honest status: T4 is a genuine
\emph{falsification} target (a T4-violating configuration kills
$B'$(i) outright), while the \emph{proof} direction needs T4 $+$ a
completion principle still to be found.
\redflag any future note claiming ``T4 $\Rightarrow$ $B'$(i)''
unconditionally.

\gap{4.6} \emph{Small remainders} (from the receipt's gap list):
C11's empty \emph{total} intersection is a limit statement (script
checks the finite shadow --- the two-class witness on $\mathbb N$ is
three lines by hand); the general-$\Omega$ halves of C1/C1$'$/C2/C3
are the written proofs in \S11a--b (each under ten lines; the machine
corroborated all finite instances); the monotonicity strengthening
(any concrete $\sigma$-class, two lines: $E \subseteq A$, $\mu(E) =
1$, $\mu(A) = 0$ gives $\mu(E \uplus A^c) = 2$).

%=====================================================================
\section{Item 5 (s7 + s10, reconstruction lane): Theorem P, Lemma NG,
Theorem B (the pruning stack)}

\noindent\textbf{Claim chain:} Theorem P (exact characterization,
all $k$, all $L$, no strong-connectivity hypothesis:
$\mathrm{TR}_k(\text{rot-1}) = \mathrm{LISC}_k$); Lemma NG
(non-generator rotations overcount: orbit decomposition to
$\mathrm{LISC}_{k/d}$); Theorem B (effective certificate: $k$-cap,
first-repeat minimality, union bookkeeping) $+$ the
$\mathrm{Safe}(\rho)$ instrument.

\noindent\textbf{Receipt:}
\file{papers/reconstruction/notes/PROOF\_READ\_2026-07-10\_pruning.md};
proof \file{papers/reconstruction/notes/pruning\_theorem\_and\_B.md};
scripts \file{papers/reconstruction/oracles/proof\_read\_2026-07-10/}.
Machine coverage is unusually strong here ($k \le 8$ exhaustive on
adversarial targets, including non-strongly-connected and
multi-component); the hand gaps are the general-$k$ arguments.

\smallskip\noindent\textbf{Lean bolster (s18):}
\file{formalization/QuerySystem/QuerySystem/PruningTheorem.lean}
(builds; receipts \texttt{[propext, Quot.sound]} --- not even choice):
\begin{itemize}[leftmargin=2em]
\item \textbf{Gap 5.1 DISCHARGED} --- Step 3 at general $k$
  (\texttt{isLISC\_of\_isTR\_one}): simplicity from layer-injectivity +
  tuple injectivity, strand boundaries and the closing step included;
  $\rho$ is an ARBITRARY relation on an ARBITRARY (not even finite)
  alphabet --- no strong connectivity can hide anywhere. Step 2
  (\texttt{isTR\_one\_of\_isLISC}) likewise; phase normalization is
  absorbed by the encoding (layers don't constrain $\rho$-arcs).
\item \textbf{Gap 5.2 DISCHARGED} --- Step 4
  (\texttt{isTR\_rot\_iff\_isTR\_one}): slots are `ZMod k`, the
  relabelling is $j \mapsto j \cdot \bar r$, and $(j + r)\bar r = j\bar
  r + 1$ is ring algebra --- the index bookkeeping has no off-by-one
  LEFT to hide. Composite: \texttt{pruning} = Theorem P, exact for
  every $L$, every $\rho$, every generator rotation.
\item \textbf{Gap 5.3 remains} (Theorem B's pigeonhole/first-repeat and
  the union bookkeeping --- rides on P; the instrument covers the
  computational content), as does Lemma NG's orbit decomposition
  (refutation half machine-checked by the oracle).
\end{itemize}

\gap{5.1} \emph{Step 3 (assembly $\Leftarrow$) at general $k$.} The
walks re-assembled from $\mathrm{TR}^{(1)}$ witnesses are simple, of
length $kL$, ring-consistent --- verified computationally on five
targets chosen to tempt collisions, but the general argument is the
written proof's.
\lookfor the simplicity argument: layer-injectivity (Lemma 0's layer
clock) is what forbids revisits; check the proof never assumes strong
connectivity (the receipt confirms the oracle cross-check included
non-SC and multi-component graphs, so the \emph{statement} is right;
your check is that the \emph{written argument} also avoids it).

\gap{5.2} \emph{Step 4 (generator conjugation) at general $k$.} The
identity $P(\sigma_1(u))[j] = \sigma_r(P(u))[j]$ reduces to $r \bar r
\equiv 1 \pmod k$; machine-checked for all 21 pairs with $k \le 8$.
\lookfor the general congruence argument is three lines of modular
arithmetic in the written proof --- verify the index bookkeeping once
(this is where an off-by-one would hide; the machine found none up to
8).

\gap{5.3} \emph{Theorem B(1)'s pigeonhole and the $k \ge 2$ scoping.}
Unsafe is a union over $2 \le k \le |A|$; the $k$-cap is pigeonhole
$+$ Theorem P.
\lookfor the union bounds match the instrument's
\texttt{range(2, A+1)} (receipt confirms); the $k = 1$ remark in the
proof is inert reader-reassurance (receipt finding 2) --- nothing to
verify beyond noticing it claims nothing.

\gap{5.4} \emph{Citations flagged non-load-bearing.} Wielandt /
Schwarz / Dulmage--Mendelsohn are bridge references only; the brute
repeat-search is what is certified.
\lookfor confirm the proof text flags them as such (it does per the
receipt); no action unless you want the classical anchors verified for
the eventual write-up. Cosmetic: instrument targets adv2 $\cong$ adv3
(relabeling), so ``15 targets'' is 14 distinct --- worth a one-word
edit only if the instrument log is ever quoted.

%=====================================================================
\section*{After ratification}

Tell me the outcome per item (ratify / ratify-with-notes / veto), and
I will: mark the taxonomy rows and MEMORY hooks RATIFIED with date;
strike the ``awaits user ratification'' banners in the attack note,
frontier map, and shovel plan; and (for any veto) open the repair as
the next session's first item. The forward queue after that is
unchanged: T4 via Maharam (8.2), coarse-block toy hunt (target: the
corrected wall statement), $B'$(ii) restatement, PP1994; parallel
relaxed-gate items FIDELITY\_REVIEW, abstract call, statistics
verdict, ILL hunt.

\end{document}
